\documentclass[letterpaper,landscape,9pt,fleqn]{extarticle}
%\usepackage[utf8]{inputenc}
%\usepackage[T1]{fontenc}
\usepackage{graphicx}
\usepackage{xcolor}
\usepackage{tikz}
\usepackage{url}
\usepackage{qrcode}
%\usetikzlibrary{shapes.geometric}
%\usetikzlibrary{calc}
%\usepackage{fourier}
\usepackage{graphicx,nicefrac}
\usepackage{isomath,upgreek,xcolor,comment}
\usepackage{pdfpages}
%\usepackage{tkz-euclide}
%\usetkzobj{all}
\pagestyle{empty}
\usepackage[activate={true,nocompatibility},final,tracking=true,kerning=true,factor=1100,stretch=10,shrink=10]{microtype}
\usepackage[american]{babel}
\usepackage{centernot}
\usepackage{amstext} % for \text macro
\usepackage{array}   % for \newcolumntype macro
\newcolumntype{L}{>{$}l<{$}} % math-mode version of "l" column type

\newcommand{\dom}{\mathrm{dom}} 
\newcommand{\range}{\mathrm{range}} 
\newcommand{\codom}{\mathrm{codom}}
\newcommand{\zero}{\mathrm{zero}} 
\newcommand{\reals}{\mathbf{R}} 
\newcommand{\ball}{\mathrm{ball}}
\newcommand{\integers}{\mathbf{Z}} 
\newcommand{\ssep}{\mid}
\newcommand{\arcsec}{\mathrm{arcsec}}
\newcommand{\arccsc}{\mathrm{arccsc}}
\newcommand{\arccot}{\mathrm{arccot}}
\newcommand{\glb}{\mathrm{glb}}
\newcommand{\lub}{\mathrm{lub}}
\newcommand{\length}{\mathrm{length}}
\usepackage{mathtools}
\DeclarePairedDelimiter{\parens}{\lparen}{\rparen}
\usepackage{amsmath,amssymb,textcomp}
\everymath{\displaystyle}

%\usepackage{times}
%\renewcommand\familydefault{\sfdefault}
%\usepackage{tgheros}
%\usepackage[defaultmono,scale=0.85]{droidmono}
%\usepackage{fourier}
\usepackage{multicol}
\setlength{\columnseprule}{0pt}
\setlength{\columnsep}{20.0pt}

\usepackage[]{enumerate}
\usepackage{expdlist}

\newenvironment{alphalist}{
  \begin{enumerate}[(a)]
    \addtolength{\itemsep}{-1.0\itemsep}}
  {\end{enumerate}}

\usepackage{geometry}
\geometry{letterpaper,left=0.4in,right=0.4in,top=0.4in,bottom=0.4in}

%\linespread{1.3}


% custom title
\makeatletter
\renewcommand*{\maketitle}{%
\noindent
\begin{minipage}{0.4\textwidth}
\begin{tikzpicture}
\node[rectangle,rounded corners=6pt,inner sep=10pt,fill=blue!50!black,text width= 0.95\textwidth] {\color{white}\Huge \@title};
\end{tikzpicture}
\end{minipage}
\hfill
\begin{minipage}{0.55\textwidth}
\begin{tikzpicture}
\node[rectangle,rounded corners=3pt,inner sep=10pt,draw=blue!50!black,text width= 0.95\textwidth] {\LARGE \@author};
\end{tikzpicture}
\end{minipage}
\bigskip\bigskip
}%
\makeatother

% custom section
\usepackage[explicit]{titlesec}
\newcommand*\sectionlabel{}
\titleformat{\section}
  {\gdef\sectionlabel{}
   \normalfont\sffamily\Large\bfseries\scshape}
  {\gdef\sectionlabel{\thesection\ }}{0pt}
  {
\noindent
\begin{tikzpicture}
\node[rectangle,rounded corners=3pt,inner sep=4pt,fill=blue!50!black,text width= 0.95\columnwidth] {\color{white}\sectionlabel#1};
\end{tikzpicture}
  }
\titlespacing*{\section}{0pt}{15pt}{10pt}


% custom footer
%\usepackage{fancyhdr}
%\makeatletter
%\pagestyle{fancy}
%\fancyhead{}
%\fancyfoot[C]{\footnotesize \textcopyright\ \@date\ \ \@author}
%\renewcommand{\headrulewidth}{0pt}
%\renewcommand{\footrulewidth}{0pt}
%\makeatother

\raggedbottom 
\raggedright
%\usepackage{tikz-3dplot}
\begin{document}

%\maketitle

\begin{multicols*}{3}

\section*{Greek characters}
\begin{tabular}{|L | L| L|} \hline
\mbox{Name} & \mbox{Symbol} & \mbox{Typical use(s)} \\ \hline
\mathrm{alpha} & \alpha  & \mbox{angle, constant} \\
\mathrm{beta} & \beta  & \mbox{angle, constant}  \\ 
\mathrm{gamma} & \gamma & \mbox{angle, constant} \\
\mathrm{delta} & \delta  & \mbox{limit definition}\\
\mathrm{epsilon} & \epsilon  \mbox{ or } \varepsilon & \mbox{limit definition} \\
\mathrm{theta}  & \theta  \mbox{ or } \vartheta &\mbox{angle}\\ 
\mathrm{pi} & \pi \mbox{ or } \uppi & \mbox{circular constant} \\
\mathrm{phi} & \phi \mbox{ or } \varphi  & \mbox{angle, constant} \\
\hline
\end{tabular}

\section*{Named sets}
\begin{minipage}[l]{0.15\textwidth}
    \begin{tabular}{|L | L |} \hline 
        \mathrm{empty\,\, set} & \varnothing \\ 
        \mathrm{real\,\, numbers} & \mathbf{R} \\
        \mathrm{ordered \, \, pairs }  & \mathbf{R}^2 \\
        \hline
    \end{tabular}
\end{minipage}
    \begin{minipage}[l]{0.15\textwidth}
        \begin{tabular}{|L | L |} \hline 
            \mathrm{integers } & \mathbf{Z} \\
            \mathrm{positive \,\,  integers } & \mathbf{Z}_{>0} \\ 
            \mathrm{positive \,\, reals} & \mathbf{R}_{>0} \\
            \hline
        \end{tabular}              
\end{minipage}

\section*{Set symbols}
\begin{minipage}[l]{0.15\textwidth}
    \begin{tabular}{|L | L|} \hline
        \mbox{Meaning}  & \mbox{Symbol} \\ \hline
        \mathrm{is \,\, a \,\, member} & \in \\
        \mathrm{subset}       & \subset \\
        \mathrm{intersection} & \cap \\ \hline
        \end{tabular}   
\end{minipage}
\begin{minipage}[l]{0.15\textwidth}
    \begin{tabular}{|L | L|} \hline
        \mbox{Meaning}  & \mbox{Symbol} \\ \hline
        \mathrm{union} & \cup  \\ 
        \mathrm{complement} & \mbox{superscript}^\mathrm{C} \\
        \mathrm{set \,\,  minus}  & \setminus \\ \hline
        \end{tabular}   
\end{minipage}

\section*{Logic symbols}
\begin{minipage}[l]{0.15\textwidth}
    \begin{tabular}{|L | L|} \hline 
        \mbox{Meaning}  & \mbox{Symbol} \\ \hline 
        \mathrm{negation} &  \lnot   \\
        \mathrm{and} &  \land  \\
        \mathrm{or} &  \lor  \\
        \mathrm{implies} &  \implies \\ \hline    
    \end{tabular}   
\end{minipage}
\begin{minipage}[c]{0.15\textwidth}
    \begin{tabular}{|L | L|} \hline 
        \mbox{Meaning}  & \mbox{Symbol} \\ \hline 
        \mbox{equivalent} &  \equiv \\ 
        \mbox{iff} & \iff \\ 
        \mbox{for all} & \forall \\
        \mbox{there exists} & \exists \\ \hline
    \end{tabular}   
\end{minipage}

\begin{minipage}[t]{0.3333\textwidth}
\section*{Truth Tables}
    \begin{tabular}{|L | L| L | L| L | L|} \hline 
        P & Q & P \land Q & P \lor Q & P \implies Q &  P \equiv Q \\ \hline \hline 
        T & T & T & T & T & T \\
        T & F & F & T & T & F \\
        F & T & F & T & T & F \\
        F &  F & F & F & T & T \\ \hline
    \end{tabular}
\end{minipage}

\begin{minipage}[t]{0.3333\textwidth}
\section*{Tautologies}
\vspace{-0.1in}
\begin{alignat*}{1}
    & \lnot \lnot P \equiv P \\
    & (P \lor P) \equiv P \\
    & (P \land  P) \equiv P \\
    & (P \equiv Q) \equiv (Q \equiv P) \\
    &\parens*{P \implies Q} \equiv \parens*{ P \lor \lnot Q}\\
    &\parens*{P \not \implies Q} \equiv \parens*{ \lnot P \\land  Q}\\
    &\lnot \parens*{P \land Q} \equiv \parens*{\lnot P \lor \lnot Q} \\
    & \parens*{P \implies Q} \equiv \parens*{\lnot Q \implies \lnot P}\\
    & \parens*{P \centernot \implies Q} \equiv \parens*{P \land \lnot Q} \\
    &\parens*{P \iff Q} \equiv \parens*{\parens*{P \implies Q} \land \parens*{Q \implies P}} \\
    &\lnot \parens*{\forall \, x \in A}\parens*{P(x)} \equiv \parens*{\exists \, x \in A}\parens*{\lnot P \parens*{x}}  \\
    &\lnot \parens*{\exists \, x \in A} \parens*{P \parens*{x}} \equiv \parens*{\forall \, x \in A} \parens*{\lnot P \parens{x}}
 \end{alignat*}
\end{minipage}

\section*{Arithmetic properties}
\begin{alignat*}{2}
    &\parens*{\forall a,b\in \reals} \parens*{a + b = b + a} & \mbox{commutivity}\\
    &\parens*{\forall a,b,c\in \reals} \parens*{a + (b + c)  = (a+b)+c} & \mbox{commutivity}\\
    &\parens*{\forall a,b \in \reals} \parens*{a  b = b a} & \mbox{commutivity}\\
    &\parens*{\forall a,b,c \in \reals} \parens*{a  (b c) = (a b) c} & \mbox{commutivity}\\
    &\parens*{\forall a,b,c \in \reals} \parens*{a (b+c) =a b + a c} & \mbox{distrutity}\\
\end{alignat*}
\section*{Function notation}
\begin{tabular}{|L | L|} \hline 
    \dom(F) &   \mbox{domain of function } F \\
    \range(F) &   \mbox{range of function } F \\
    \mathrm{C}_{A} & \mbox{ set of continuous functions on set } A \\
    \mathrm{C}_{A}^1 & \mbox{ set of differentiable functions on set } A \\
    A \to B   & \mbox{set of functions from } A \mbox { to } B \\ \hline
\end{tabular}

\section*{Set operators}
\begin{alignat*}{1}
    A \cup B &= \{ x \mid x \in A \lor  x \in B \} \\
    A \cap B &= \{ x \mid x \in A \land  x \in B \} \\
    A \setminus B &= \{ x \mid x \in A \land  x \notin B \} \\
    A \times B &= \{ (a,b) \mid a \in A \land  b \in B \}
\end{alignat*}

 \section*{Generalized set operators}
 Each member of a set $\mathcal{C}$ is a set:
\begin{alignat*}{1}
    &\underset{A \in \mathcal{C}}{\bigcup} A = \left \{z \mid \parens*{\exists \, B  \in \mathcal{C} }\parens{z \in B} \right \}\\
    &\underset{A \in \mathcal{C}}{\bigcap} A = \left \{z \mid \parens*{\forall \, B \in \mathcal{C}}\parens{z \in B} \right \}
\end{alignat*}
Theorem: \(\underset{A \in \mathcal{C}}{\bigcup} A^\mathrm{C} = \parens*{\underset{A \in \mathcal{C}}{\bigcap} A}^\mathrm{C} \)
\section*{Functions applied to  sets}
Let $A \subset \dom(F)$ and $B \subset \range(F)$:
\begin{alignat*}{1}
    F(A) &= \{F(x) \mid x \in A \} \\
    F^{-1}(B) &= \{x \in \dom(F) \mid F(x) \in B \}
\end{alignat*}


\section*{Triangle inequalities}
For all $x,y \in \reals$, we have
\begin{alignat*}{1}
    |x+y| &\leq |x| + |y| \\
    \big | |x| - |y|  \big |  &\leq |x-y|    
\end{alignat*}
\section*{Floor and ceiling}
\vspace{-0.1in}
\noindent Definitions:
\begin{align*}
    \lfloor x \rfloor = \max \{k \in \integers \mid  k \leq x \} \\
    \lceil x \rceil = \min  \{k \in \integers \mid  k \geq x \}   
\end{align*}

\noindent Properties:
\begin{align*}
  % &\parens*{\forall x \in \reals} \parens*{\lfloor x \rfloor \leq x} \\
  % &\parens*{\forall x \in \reals} \parens*{\lceil x \rceil \geq x} \\
   & \parens*{\forall x \in \reals, n \in \integers} \parens*{x < n \iff  \lfloor x \rfloor < n}\\
   & \parens*{\forall x \in \reals, n \in \integers} \parens*{n < x \iff  n < \lceil  x \rceil}\\
\end{align*}
\vspace{-0.51in} 


\section*{Elementary function properties}
    \begin{description}[\itemsep=0em]
        \item[Increasing] \( \parens*{\forall \, x,y \in A} \parens*{x < y \implies F(x) \leq F(y)} \).
        For strictly increasing, replace $F(x) \leq F(y)$ with $F(x) < F(y)$.
        \item[Decreasing] \( \parens*{\forall \, x,y \in A} \parens*{x < y \implies F(x) \geq F(y)} \)
        For strictly decreasing, replace $F(x) \geq F(y)$ with $F(x) > F(y)$.
        \item[One-to-one] \( \parens*{\forall \, x,y \in \dom(F)} \parens*{F(x) = F(y) \implies x = y} \)
    \end{description}

\section*{Equivalence relations}
Let $R \in A \times A \to \{\mbox{true}, \mbox{false}\}$. We say
\begin{description}
    \item[reflective]  $\parens*{\forall x \in S} \parens*{x \, R \, x}$
    \item[symmetric] $\parens*{\forall x,y\in S} x \, R \, y \implies y \, R \, x$
    \item[transitive] $\parens*{\forall x,y,z\in S} x \, R \, y \land y \, R \, z
       \implies x \, R \, z$
\item[Equivalence class] $[x] = \{s \in S \mid s \, R \, x \}$
\end{description}

\section*{Axioms}
\begin{description}
      \item[Well-ordering]  Every nonempty set of positive integers 
        contains a least element.
    \item[Induction] $\parens*{\forall n \in \integers_{\geq 0}}(P(n))$ if and only if
    $ P(0) \land \parens*{\forall n \in 
    \integers_{\geq 0}}  \parens*{P(n) \implies P(n+1)}$.
\end{description}


\vfill
\noindent {Revised \today. Barton Willis is the author of this work. This work is
licensed under Attribution 4.0 International (CC BY 4.0) \,  \qrcode[height=0.15in]{https://creativecommons.org/licenses/by/4.0/}. For the current version of
this document, visit \, \qrcode[height=0.15in]{https://github.com/barton-willis/Math-100-200-level}}
\end{multicols*}%{3}

\end{document}